\documentclass[a4paper,10pt,fleqn]{article}
\usepackage[margin=1in]{geometry}
\usepackage{amssymb}
\usepackage{adjustbox}
\usepackage{natbib}
\usepackage[colorlinks=true,linkcolor=blue,citecolor=blue,urlcolor=blue]{hyperref}
\usepackage{fuzz}
\usepackage{schemapk}
\usepackage{zed-maths}
\usepackage{zed-proof}
\newdimen\savedleftskip
\begin{document}

\section*{Airline Booking State}

\noindent Define the airline state schema, which holds the set of routes and the bookings relation.

\bigskip

\begin{zed}
[Airline, BookingId, CustomerId, RouteId]
\end{zed}

\begin{schema}{AirlineState}
routes : \power RouteId \\
bookings : BookingId \pfun CustomerId \cross RouteId
\end{schema}

\noindent A read-only operation that queries whether a route exists uses Xi to signal that the state is unchanged.

\bigskip

\begin{schema}{RouteExists}
\Xi AirlineState \\
routeId? : RouteId \\
result! : \nat
\where
routeId? \in routes \\
result! = 1
\end{schema}

\noindent An operation that adds a booking uses Delta to include both the before- state and the after-state. The primes on bookings' and routes' are Z convention for the output state; here bookings' is added as a typed declaration because the assignment is the predicate.

\bigskip

\begin{schema}{AddBooking}
\Delta AirlineState \\
bookingId? : BookingId \\
customerId? : CustomerId \\
routeId? : RouteId
\where
bookingId? \notin \dom bookings \\
routeId? \in routes
\end{schema}

\section*{Stack with Generic Inclusion}

\noindent Schema inclusion works with generic types. Delta Stack[N] brings in the before-and-after state of a stack instantiated over the naturals.

\bigskip

\begin{schema}{StackOp}[X]
delta : \nat
\end{schema}

\begin{schema}{PushNat}
\Delta StackOp[\nat] \\
value? : \nat
\where
delta > 0
\end{schema}

\end{document}